% No modificar estas líneas de código, por favor dirigirse a INTRODUCCIÓN 
\newpage
\phantomsection\addcontentsline{toc}{section}{INTRODUCCIÓN} 

%------------------------
%
%       INTRODUCCIÓN
%
%------------------------

\section*{INTRODUCCIÓN}

En los procesos de manufactura moderna, el mecanizado por arranque de viruta mediante tornos de Control Numérico Computarizado (CNC) representa una de las operaciones más críticas y demandadas para la fabricación de componentes de alta precisión. Dentro de este proceso, la herramienta de corte es el elemento primario que interactúa directamente con el material de trabajo, estando sometida continuamente a severos esfuerzos mecánicos, elevadas temperaturas y fricción dinámica. Como consecuencia natural de esta interacción termomecánica, los insertos experimentan una degradación progresiva en la geometría de su filo.

La evolución del desgaste de las herramientas incide de forma determinante en la estabilidad del proceso de corte, modificando las fuerzas de mecanizado y alterando directamente la calidad del acabado superficial y la precisión dimensional de las piezas manufacturadas. Tradicionalmente, la gestión del reemplazo de insertos en la industria se ha sustentado en inspecciones visuales periódicas o en tablas estandarizadas de vida útil. Sin embargo, estas estrategias empíricas resultan inadecuadas para anticipar fallas dinámicas o desgastes anómalos generados por variaciones en las propiedades del material o fluctuaciones en las condiciones de corte.

Frente a estas limitaciones operativas, el desarrollo de sistemas de Monitoreo de Condición de Herramientas (\textit{Tool Condition Monitoring}, TCM) se ha consolidado como una disciplina fundamental dentro de la manufactura inteligente. El monitoreo indirecto y continuo permite inferir en tiempo real el estado de la herramienta a través de señales físicas del proceso, como las variaciones en la corriente del motor del husillo y las oscilaciones mecánicas de alta frecuencia, facilitando una transición hacia estrategias avanzadas de mantenimiento predictivo.

En este contexto tecnológico, el presente proyecto de grado tiene como propósito diseñar e implementar un sistema periférico no invasivo de monitoreo de condición para tornos CNC, integrando adquisición de señales de corriente y vibración, procesamiento inteligente en el borde (\textit{Edge AI}) y conectividad industrial. La investigación y validación experimental del sistema se desarrollan específicamente en el entorno operativo de la empresa metalmecánica Maestranza San Carlos (MAINSA), adaptando la solución a las exigencias reales de su línea de producción.

\subsection*{Antecedentes}
\phantomsection\addcontentsline{toc}{subsection}{Antecedentes}

El monitoreo de condición de herramientas de corte constituye una estrategia relevante para mejorar la calidad de manufactura. Permite disminuir la generación de piezas defectuosas y reducir paradas no planificadas en máquinas CNC. Dentro de las técnicas indirectas de monitoreo, el análisis de señales eléctricas de los accionamientos representa una alternativa de bajo impacto. Esto permite inferir cambios en la carga de corte sin necesidad de instalar sensores invasivos directamente sobre la herramienta de corte.

En este enfoque, las variaciones de desgaste, pérdida de filo o rotura modifican las fuerzas de mecanizado. En consecuencia, se altera la corriente consumida por el motor del husillo o por los motores de avance. Un antecedente temprano y fundamental de este enfoque es el trabajo de \cite{romero2003}, quienes propusieron detectar la rotura de herramientas en fresadoras CNC. Lograron esto mediante el análisis continuo de la corriente de los accionamientos durante el mecanizado.

Su propuesta permitió identificar eventos asociados a la rotura sin recurrir a dinamómetros ni a sensores de fuerza de alto costo. Además, evitó realizar modificaciones importantes en la estructura mecánica de la máquina, lo que consolidó el uso de señales internas. A este enfoque indirecto y pragmático se le denomina frecuentemente en la literatura como monitoreo \textit{sensorless} o de instrumentación mínima.

La línea de investigación fue ampliada por \cite{trejo2010}, quien estudió el monitoreo y modelado fenomenológico del desgaste de herramienta. Las pruebas se realizaron bajo condiciones variables de corte en torno CNC, integrando señales de corriente y vibración. Se utilizó un sensor inteligente basado en FPGA con el propósito fundamental de estimar cuantitativamente la evolución del desgaste de flanco a lo largo del proceso.

Este antecedente es importante porque no se limita a identificar una falla súbita aislada durante la operación. Al contrario, aborda la evolución progresiva del desgaste y su comportamiento en condiciones de corte dinámicas y variables. Posteriormente, \cite{sevilla2011} validaron un método de detección de rotura de herramienta en fresado de alta velocidad analizando las señales de corriente del motor de avance.

El estudio de \cite{sevilla2011} utilizó la transformada wavelet discreta para extraer información transitoria de la señal analizada. Además, aplicó herramientas estadísticas robustas para discriminar eficazmente condiciones normales y anormales de mecanizado. Este valioso aporte respalda que las señales de corriente pueden contener patrones sumamente útiles para la detección de eventos súbitos aplicando procesamiento tiempo-frecuencia.

Como complemento alternativo, \cite{ramirez2018} desarrollaron un sensor inteligente para detección de rotura en fresado utilizando termografía infrarroja. Estas pruebas experimentales se realizaron rigurosamente bajo condiciones de corte seco y también utilizando abundante refrigerante. Sin embargo, la termografía requiere forzosamente una línea de visibilidad directa hacia la zona de corte en el interior de la máquina.

Este requerimiento visual puede verse afectado negativamente por la presencia de viruta, refrigerante o diversas restricciones físicas de la máquina CNC. Por ello, el monitoreo basado en el análisis de corriente mantiene una ventaja sumamente relevante y pragmática frente a los sistemas ópticos. Destaca de manera especial cuando se busca una implementación robusta de fácil instalación y de muy bajo costo computacional.

La vigencia del monitoreo no invasivo por corriente se confirma fehacientemente con investigaciones recientes que incorporan aprendizaje profundo. \cite{lai2025} propusieron un sistema eficiente basado en señales de corriente del husillo principal. Dichas señales fueron obtenidas con un sensor no invasivo de tipo SCT013 instalado externamente alrededor del cableado, sin intervenir para nada el controlador.

Los autores analizaron las señales mediante la transformada rápida de Fourier y compararon diversos modelos de redes neuronales artificiales (ANN, DNN, CNN). Su objetivo primordial fue clasificar eventos de rotura, logrando detectar graves anomalías con una extraordinaria precisión y destacable rapidez. La utilización de una pinza SCT013 es particularmente pertinente para instrumentación de muy bajo impacto y fácil despliegue.

El sensor SCT013 permite implementar un sistema de adquisición de datos totalmente reversible y seguro para maquinaria heredada sin conectividad nativa. No obstante, el diseño meticuloso del sistema debe considerar rigurosamente el rango nominal de la corriente de trabajo y la frecuencia de muestreo adecuada. También requiere implementar de manera obligatoria un óptimo acondicionamiento de señal para lograr rechazar el elevado ruido eléctrico industrial en planta.

Por otra parte, \cite{wang2024} propusieron una innovadora metodología de predicción basada estrictamente en la fusión de información multisensorial. El estudio científico integra conjuntamente complejas señales de corriente, niveles de vibración mecánica y medición directa de la fuerza de corte. Esto demuestra fehacientemente que las variables eléctricas y mecánicas siempre logran aportar información complementaria para un diagnóstico holístico.

Una limitación frecuente de los modelos tradicionales es su drástica disminución de desempeño general cuando cambian drásticamente las condiciones operativas de trabajo. \cite{huang2026} abordaron y resolvieron este problema directamente mediante la estratégica fusión de señales de vibración y corriente. Esta técnica se complementó excepcionalmente con un esquema avanzado de algoritmos de aprendizaje por transferencia completamente supervisado.

Su estrategia combina de manera sumamente eficiente el preentrenamiento algorítmico y un riguroso ajuste inicial de parámetros. Aunque la metodología requiere obtener algunos datos etiquetados nuevos, minimiza de forma muy importante la valiosa información necesaria para lograr un reentrenamiento efectivo. La calidad del procesamiento final de las señales constituye otro aspecto fundamental, crítico y determinante en la fiabilidad.

Las señales medidas operativamente en una máquina CNC pueden contener elevado ruido eléctrico, indeseables vibraciones estructurales y cambios de carga sumamente severos. Por esta precisa razón, resulta estrictamente necesario implementar siempre procedimientos eficaces de filtrado digital y segmentación temporal avanzada. En este ámbito técnico, \cite{jia2025} propusieron un método de reducción de ruido basado eficientemente en descomposición modal variacional adaptativa.

De una manera bastante similar, el detallado estudio sobre la técnica ICFO-SVMD propuso una prometedora metodología que integra optimización algorítmica y descomposición modal \citep{icfo2026}. El principal objetivo analítico fue lograr procesar señales marcadamente no estacionarias que estaban severamente afectadas por ruido electromagnético de origen industrial. En consecuencia, el presente proyecto aplicará un esquema de procesamiento robusto y fuertemente escalonado por etapas.

Además de requerir una constante mejora algorítmica, la incorporación de novedosa conectividad IoT logra potenciar sustancialmente las actuales capacidades del monitoreo. Permite fundamentalmente que el sistema embebido se convierta rápidamente en una moderna plataforma remota de apoyo a la toma de decisiones. El interesante caso de estudio aplicado de la prestigiosa Universidad de Sheffield se erige como un antecedente sumamente relevante \citep{sheffield2021}.

Finalmente, \cite{hamasha2025} propusieron un complejo y completo marco integral tecnológico e industrial de indudable y gran envergadura. Dicho marco funcional articula eficientemente diversos sensores IoT, procesamiento local en el extremo o borde operativo (\textit{Edge AI}), analítica predictiva y monitoreo remoto. El estudio logra identificar cuantiosos beneficios clave como la reducción drástica de latencia y tiempos de respuesta, advirtiendo enormes desafíos en ciberseguridad industrial.

A partir de todos los numerosos antecedentes exhaustivamente revisados, se observa que ya existen sólidas investigaciones consolidadas relativas a detección y fallas. Sin embargo, dichos estudios literarios abordan los diferentes componentes tecnológicos mayoritariamente de manera totalmente independiente y sumamente aislada. Por consiguiente, el gran y ambicioso aporte ingenieril de este proyecto radica en la inusual integración simultánea de estas tecnologías operando en el entorno real de MAINSA.

\subsection*{Identificación del problema}%
\phantomsection\addcontentsline{toc}{subsection}{Identificación del problema}
La problemática central en la empresa metalmecánica Maestranza San Carlos (MAINSA) radica en la carencia de un sistema periférico y no invasivo para el monitoreo en tiempo real del desgaste de las herramientas de corte en sus tornos CNC industriales. Actualmente, la supervisión de los insertos depende por completo de la inspección visual periódica y de la experiencia subjetiva del operador, una práctica tardía y discontinua que resulta incapaz de advertir anomalías incipientes en el filo antes de que se manifieste una falla severa.

El origen técnico de esta situación se sustenta en tres causas fundamentales identificadas en planta: primero, las máquinas herramientas no cuentan con sensores ni módulos de adquisición continua para registrar las condiciones de corte; segundo, el diagnóstico humano visual no permite cuantificar la evolución del desgaste en plena operación; y tercero, la arquitectura cerrada de los controladores CNC bloquea el acceso a las señales internas de los servoamplificadores, impidiendo implementar monitoreo por corriente sin vulnerar la integridad o garantía de los equipos.

Esta problemática genera efectos operativos y económicos altamente perjudiciales para MAINSA: la degradación no advertida del filo produce piezas fuera de tolerancia dimensional y con rugosidad superficial deficiente, demandando costosos reprocesos y reposición de material; la rotura imprevista de herramientas provoca paradas no planificados que alteran la programación de entrega; y los operarios sustituyen prematuramente los insertos ante la incertidumbre de su estado, desperdiciando vida útil remanente y encareciendo sensiblemente el costo por consumibles.

\subsection*{Formulación del problema}%
\phantomsection\addcontentsline{toc}{subsection}{Formulación del problema}
A partir del diagnóstico técnico y la problemática operacional fundamentada en el árbol de problemas de la empresa MAINSA, se establece la interrogante central que orienta la presente investigación:

¿Cómo diseñar e implementar un sistema periférico no invasivo para el monitoreo en tiempo real del desgaste de herramientas de corte en los tornos CNC de la empresa MAINSA, que permita superar la dependencia de la inspección visual subjetiva, prevenir paradas no planificadas por rotura de filo y reducir las mermas operativas en el taller?

\subsection*{Objetivos}%
\phantomsection\addcontentsline{toc}{subsection}{Objetivos}

\subsubsection*{Objetivo general}%
\phantomsection\addcontentsline{toc}{subsubsection}{Objetivo general}
Diseñar e implementar un sistema periférico no invasivo para el monitoreo en tiempo real del desgaste de herramientas de corte en tornos CNC de la empresa MAINSA, mediante la medición de variables de corriente y vibración procesadas con algoritmos de inteligencia artificial en el borde (\textit{Edge AI}), para anticipar oportunamente fallas críticas en los insertos y reducir paradas no planificadas en el proceso de mecanizado.

\subsubsection*{Objetivos específicos}%
\phantomsection\addcontentsline{toc}{subsubsection}{Objetivos específicos}
\begin{itemize}
\item Recopilar y analizar la información técnica y operativa sobre el proceso de mecanizado y los mecanismos de desgaste de herramientas en tornos CNC de MAINSA.
\item Diseñar la arquitectura mecatrónica del sistema periférico de monitoreo, integrando acondicionamiento de señal, procesamiento en el borde y conectividad IoT.
\item Seleccionar y dimensionar los sensores de corriente, vibración y la unidad de procesamiento embebido.
\item Implementar e integrar el sistema periférico de adquisición, procesamiento inteligente y supervisión remota.
\item Validar experimentalmente el funcionamiento del sistema integrado mediante pruebas de corte en los tornos CNC de MAINSA.
\end{itemize}

\subsubsection*{Acciones por objetivo específico}%
\phantomsection\addcontentsline{toc}{subsubsection}{Acciones por objetivo específico}
Con el propósito de operativizar cada uno de los objetivos específicos y detallar la metodología de ejecución técnica, en la \tab{acciones_objetivos} se desglosan las acciones concretas asociadas a cada fase del proyecto.

\begin{longtable}{|L{4.8cm}|L{9.2cm}|}
\caption{Acciones asociadas a los objetivos específicos}
\label{tab:acciones_objetivos}
\\
\hline
\textbf{Objetivo específico} & \textbf{Acciones a desarrollar} \\
\hline \hline
\endfirsthead

\hline
\textbf{Objetivo específico} & \textbf{Acciones a desarrollar} \\
\hline \hline
\endhead

\hline
\endfoot

\multicolumn{2}{c}{\textbf{Fuente:} Elaboración propia}
\endlastfoot

1. Recopilar y analizar información técnica y operativa sobre mecanizado y desgaste de herramientas en MAINSA. & 
\textbullet\ Revisión técnica de manuales de operación, cinemática y parámetros de corte de los tornos CNC en planta.\\
& \textbullet\ Caracterización de las geometrías de insertos de carburo y tipos de aceros mecanizados habitualmente.\\
& \textbullet\ Análisis de criterios y patrones de desgaste de flanco según la norma internacional ISO 3685.\\
& \textbullet\ Relevamiento de datos históricos sobre paradas no programadas, defectos superficiales y roturas de filo.\\ \hline

2. Diseñar la arquitectura mecatrónica del sistema periférico de monitoreo, acondicionamiento, Edge AI e IoT. & 
\textbullet\ Definición del diagrama de bloques funcional (adquisición, acondicionamiento, procesamiento y enlace IoT).\\
& \textbullet\ Diseño del circuito de acondicionamiento para rectificación y conversión a niveles proporcionales al valor eficaz (RMS).\\
& \textbullet\ Diseño esquemático y ruteo de la placa de circuito impreso (PCB) con aislamiento galvánico industrial.\\
& \textbullet\ Concepción del gabinete protector y soportes de fijación no invasiva para la instalación en la máquina.\\ \hline

3. Seleccionar y dimensionar los sensores de corriente, vibración y la unidad de procesamiento embebido. & 
\textbullet\ Dimensionamiento del rango de medición y relación de transformación del transductor de corriente según la potencia del motor.\\
& \textbullet\ Selección del acelerómetro considerando sensibilidad, ancho de banda y rango dinámico de vibración en la torreta.\\
& \textbullet\ Selección del microcontrolador evaluando periféricos ADC, DMA, instrucciones DSP y memoria requerida.\\
& \textbullet\ Cálculo y selección de componentes pasivos de filtrado analógico, resistencias de carga y fuentes de alimentación.\\ \hline

4. Implementar e integrar el sistema periférico de adquisición, procesamiento inteligente y supervisión remota. & 
\textbullet\ Fabricación, ensamble y verificación eléctrica de la tarjeta de circuito impreso y módulos acondicionadores.\\
& \textbullet\ Programación del firmware para muestreo síncrono, filtrado digital y extracción de descriptores estadísticos temporales.\\
& \textbullet\ Entrenamiento, cuantización y despliegue del modelo de aprendizaje automático ligero para clasificación de desgaste.\\
& \textbullet\ Configuración del protocolo de comunicaciones y desarrollo del panel de supervisión remota para visualización y alertas.\\ \hline

5. Validar experimentalmente el funcionamiento del sistema integrado mediante pruebas de corte en MAINSA. & 
\textbullet\ Instalación y calibración del sistema periférico en el torno CNC bajo condiciones operativas de producción real.\\
& \textbullet\ Ensayos de torneado cilíndrico registrando señales con insertos en estado nuevo, desgaste moderado y crítico.\\
& \textbullet\ Medición metrológica de la rugosidad superficial ($Ra$) en probetas y verificación de tolerancias dimensionales.\\
& \textbullet\ Evaluación cuantitativa de la exactitud de clasificación y cuantificación de la reducción de paradas y mermas.\\ \hline

\end{longtable}

\subsection*{Hipótesis}
\phantomsection\addcontentsline{toc}{subsection}{Hipótesis}

La integración de un sistema embebido de monitoreo no invasivo basado en la adquisición y fusión de señales de corriente y vibración mecánica, procesadas localmente mediante algoritmos de Inteligencia Artificial en el borde (\textit{Edge AI}) sobre una unidad de procesamiento embebido, permite clasificar el nivel de desgaste de los insertos y anticipar eventos de rotura súbita con una exactitud superior al 90\,\% en condiciones de corte reales, reduciendo significativamente las mermas de material y los tiempos improductivos por retrabajo en los tornos CNC de la empresa MAINSA.

\subsection*{Justificación}
\phantomsection\addcontentsline{toc}{subsection}{Justificación}

\subsubsection*{Justificación técnica}
\phantomsection\addcontentsline{toc}{subsubsection}{Justificación técnica}

Desde la perspectiva técnica, el parque de maquinaria CNC en talleres metalmecánicos como MAINSA suele presentar una notable heterogeneidad tecnológica, lo que dificulta la incorporación de sistemas comerciales de monitoreo propietarios de alto costo. En este contexto, el desarrollo de una arquitectura de instrumentación externa y no invasiva (\textit{retrofitting}) constituye una solución altamente viable que evita modificar la integridad del controlador numérico. La combinación complementaria de señales de corriente del husillo y acelerometría mecánica permite capturar simultáneamente los esfuerzos globales de corte y las oscilaciones dinámicas de alta frecuencia asociadas a la pérdida de filo.

Asimismo, la implementación de algoritmos de aprendizaje automático directamente en el extremo de la red (\textit{Edge AI}) sobre la unidad de procesamiento embebido aporta una ventaja operativa determinante al eliminar la latencia de transmisión hacia servidores externos. Esta capacidad de inferencia local garantiza tiempos de respuesta casi instantáneos ante anomalías severas en el mecanizado, asegurando una supervisión continua y robusta frente al elevado ruido electromagnético y las complejas perturbaciones que caracterizan el entorno productivo real del taller industrial.

\subsubsection*{Justificación económica}
\phantomsection\addcontentsline{toc}{subsubsection}{Justificación económica}

En el ámbito económico, la rotura intempestiva de insertos y el mecanizado con herramientas degradadas generan pérdidas financieras directas debido al desperdicio de materia prima y a la necesidad imperiosa de ejecutar costosos retrabajos sobre piezas fuera de tolerancia. La implementación de un sistema de mantenimiento predictivo de bajo costo permitirá maximizar la vida útil efectiva de cada inserto sin comprometer los acabados superficiales, reduciendo notablemente las paradas no programadas de máquina y optimizando el rendimiento operativo global de la empresa MAINSA.

\subsubsection*{Justificación social y operativa}
\phantomsection\addcontentsline{toc}{subsubsection}{Justificación social y operativa}

En el plano social y de seguridad industrial, la fractura violenta de una herramienta de corte durante operaciones de desbaste o cilindrado a alta velocidad representa un riesgo latente para la integridad física de los operarios debido a la posible proyección de esquirlas metálicas y potenciales colisiones mecánicas severas. Por consiguiente, dotar al personal técnico de herramientas tecnológicas que anticipen fallas críticas promueve un entorno de trabajo más seguro y confiable, fomentando a la vez la apropiación y transferencia de tecnologías de la Industria 4.0 en el sector productivo regional.

\subsection*{Alcance}
\phantomsection\addcontentsline{toc}{subsection}{Alcance}

El alcance del presente trabajo de grado comprende el diseño, desarrollo, calibración y validación experimental de un prototipo embebido funcional para el monitoreo de condición de herramientas en tornos CNC de la empresa MAINSA. La solución abarca desde la adquisición y acondicionamiento de señales analógicas de corriente (mediante transductor no invasivo) y vibración mecánica (mediante acelerometría MEMS), pasando por la ejecución del modelo de inferencia \textit{Edge AI} en la unidad de procesamiento embebido, hasta la transmisión inalámbrica de datos hacia una plataforma IoT para visualización remota en tiempo real.

\subsection*{Límites}
\phantomsection\addcontentsline{toc}{subsection}{Límites}

El proyecto se encuentra delimitado a las condiciones operativas y tipos de mecanizado habituales en el taller de MAINSA, focalizándose principalmente en operaciones de torneado sobre aceros al carbono y aleados bajo parámetros de corte estándar. Asimismo, el sistema operará como una plataforma de supervisión y diagnóstico predictivo con generación de alertas tempranas, sin intervenir directamente los circuitos de paro de emergencia (\textit{E-Stop}) ni alterar los lazos de control internos del CNC, garantizando en todo momento la preservación de las garantías de los equipos.

Desde la perspectiva del análisis eléctrico, el proyecto delimita la medición al comportamiento de la componente fundamental de corriente y voltaje alterno, asumiendo una señal senoidal pura. Por consiguiente, se excluye deliberadamente el análisis de distorsión armónica total (THD), componentes armónicos de alta frecuencia y fenómenos transitorios derivados de la modulación por ancho de pulsos (PWM) del variador de frecuencia. Las variaciones en el consumo energético se atribuirán y justificarán estrictamente a los cambios en la magnitud eficaz (RMS) de la tensión y la corriente como respuesta a la velocidad del motor y a la carga mecánica de corte, facilitando la conversión a niveles continuos (DC) para un procesamiento simplificado y eficiente en el microcontrolador.

\subsection*{Estructura del documento}
\phantomsection\addcontentsline{toc}{subsection}{Estructura del documento}

El presente documento se encuentra estructurado en acápites y capítulos correlativos conforme a las directrices metodológicas institucionales. La sección de Introducción expone la contextualización, los antecedentes, el problema, los objetivos y la justificación. El Capítulo 1 desarrolla el Marco Teórico con los fundamentos de mecanizado, señales, \textit{Edge AI} e IoT. Finalmente, los capítulos sucesivos abordan el Marco Metodológico, el diseño e implementación del prototipo práctico, concluyendo con el análisis de resultados y recomendaciones.